\section{从 AI Bubble 到 AI War：为什么“泡沫”不是充分解释}

本节先处理全片的总框架。节目开头问的是“AI Bubble”，广密给出的回答却是“更像 AI War”。这不是文字游戏，而是分析单位的变化。如果只把 AI 当成商业项目，就会从现金流、估值和资本开支出发，问“这笔钱怎样赚回来”；如果把 AI 当成下一代平台、国防能力和科技秩序重组，就还要问“如果不投，会不会被颠覆”。

\begin{figure}[H]
\centering
\includegraphics[width=0.96\textwidth]{ai-war-alliances.png}
\caption{AI War 的两条主线：Nvidia GPU 生态、Google TPU 生态，以及围绕 OpenAI/Google 形成的反制联盟。自制概念图，依据 00:15:04--00:17:48 对谈内容整理。}
\end{figure}

\begin{knowledgebox}{读图：这张图先看生态，再看联盟}
左上和中上分别是 Nvidia GPU 与 Google TPU 两条硬件生态主线。OpenAI、Anthropic 等领先模型公司更多站在 GPU 生态中；Google 则在模型、芯片、云和产品端到端整合。下方的“反 Google 联盟”和“反 OpenAI 联盟”提醒我们：越强的单一玩家，越会触发其他巨头、资本和生态伙伴的制衡。
\end{knowledgebox}

\subsection{AI War 的第一层含义：商业防守}

前面说 AI War 不是修辞，本节先看商业层面的输不起。广密举了两个现实例子。第一个是云厂商被模型白牌化：过去开发者和企业优先考虑 AWS、Azure 或 Google Cloud，现在很多新应用先问“用哪个模型”：GPT、Claude 还是 Gemini。云仍然重要，但用户心智可能从云迁移到模型。

第二个是 Agent 可能成为新流量入口。传统搜索、Super App、OTA、出行、招聘和电商平台，许多商业模式都建立在信息分发、匹配和交易抽成上。如果 Agent 能替用户比较、谈判、下单、支付和跟踪，平台的抽成和入口权就会被重新定价。对巨头来说，这不是“多一个工具”，而是已有利润池的防守战。

\begin{warningbox}{误读警告：AI War 不等于短期账一定成立}
说 AI 是战争，不代表每个数据中心项目、每个模型训练计划和每家创业公司都值得投资。它只说明关键玩家的行为不能用普通 SaaS 投资回报周期解释，因为“不投入”的机会成本可能是平台地位被重写。
\end{warningbox}

\subsection{AI War 的第二层含义：国家战略和算力动员}

广密进一步把 AI 比作新的国防和核武器。这句话容易被听成夸张，但它指向的是国家层面的能力结构：谁掌握更强模型、更多算力、更密集研究员、更完整芯片供应链，谁就可能在软件生产、科学发现、军事辅助、产业自动化和信息分发上获得系统性优势。

这也是为什么短期现金流算不过来的项目仍会获得资本和政策支持。商业公司担心被颠覆，国家担心战略落后，资本市场又在寻找少数能吸纳巨大资金并仍有上行想象的方向。三者叠加后，AI 资本开支的节奏就很难只用“泡沫”二字解释。

\begin{importantbox}{分析 AI 投入的三层账}
\begin{tabularx}{\textwidth}{p{0.20\textwidth}p{0.35\textwidth}X}
\toprule
账本 & 关心的问题 & 对应本期判断 \\
\midrule
财务账 & 收入、毛利、现金流、折旧如何覆盖投入 & OpenAI 的 1.4T 承诺短期会被质疑。 \\
平台账 & 新入口会不会替代旧入口，谁掌握用户关系 & Agent 与模型云可能重写搜索、云和 Super App。 \\
战略账 & 如果不投入，国家和公司是否失去未来能力 & AI 被类比为国防和核武器，玩家不敢轻易离场。 \\
\bottomrule
\end{tabularx}
\end{importantbox}

\subsection{从泡沫争论转向格局分析}

泡沫争论常常问“估值是不是太高”。本期更值得保留的问题是“谁会被迫继续投入、谁能形成联盟、谁能让投入变成产品和数据闭环”。这也是后文的展开顺序：先拆 OpenAI 的收入账，再看 Nvidia/Google 阵营，然后讨论三大模型为什么交替领先，最后进入 Online Learning 这个可能改变现有格局的变量。

\subsection{本章小结}

AI War 框架把讨论从估值情绪拉回结构变量。它不否认局部泡沫，却提醒读者：当 AI 同时威胁云、搜索、Super App、白领劳动和国家竞争力时，巨头和国家会用比普通商业项目更激进的方式投入。

